\documentclass[11pt]{article}
\usepackage{amsmath,amsfonts,amssymb,amsthm}
\usepackage{graphicx}
\usepackage{algorithm}
\usepackage{algorithmic}
\usepackage{float}
\usepackage{hyperref}

\title{Local Randomized Neural Networks Methods for Interface Problems}
\author{Yunlong Li and Fei Wang}
\date{}

\begin{document}
\maketitle

\section{Introduction}
Interface problems appear in many complex physical situations such as multiphase flow and fluid-structure interactions where multiphysics coupling occurs across interfaces with potentially intricate geometries and dynamic boundaries. Traditional numerical methods often face challenges in handling the coupling conditions at these interfaces, especially when they have complex shapes or move over time.

Recently, neural network-based approaches have shown potential for solving partial differential equations (PDEs) with complex geometries, but they suffer from drawbacks including time-consuming optimization processes and the risk of getting trapped in local optima. Randomized Neural Networks (RNNs) offer an alternative approach by fixing randomly selected weights and biases in hidden layers, thereby avoiding expensive optimization solvers during training.

This paper proposes Local Randomized Neural Networks (LRNNs) method, a mesh-free approach where different RNNs are used to approximate solutions in different subdomains divided by interfaces. This method discretizes interface problems into linear systems at randomly sampled points and solves them using least-squares methods.

\section{Model Interface Problems}

\subsection{Elliptic Interface Problem}
Consider a domain $\Omega$ divided by an interface $\Gamma$ into two subdomains $\Omega_1$ and $\Omega_2$. The elliptic interface problem is defined as:
\begin{align}
-\nabla \cdot (\beta(x)\nabla u) &= f, &x \in \Omega, \\
[u] &= g_1, &x \in \Gamma, \\
[\beta(x)\nabla u \cdot n] &= g_2, &x \in \Gamma, \\
u &= g_D, &x \in \partial\Omega
\end{align}
where $\beta(x)$ is a piecewise positive constant function:
\begin{equation}
\beta(x) = 
\begin{cases}
\beta_1, & x \in \Omega_1, \\
\beta_2, & x \in \Omega_2
\end{cases}
\end{equation}

The notation $[u]$ denotes the jump of $u$ across the interface, $[u] = u_1 - u_2$, and $n$ is the unit outward normal vector on $\Gamma$ from $\Omega_1$ to $\Omega_2$.

\subsection{Parabolic Interface Problem}
For time-dependent problems over interval $(0,T)$, the parabolic interface problem is:
\begin{align}
\frac{\partial u}{\partial t} - \nabla \cdot (\beta(x)\nabla u) &= f, &(x,t) \in \Omega \times (0,T), \\
[u] &= g_1, &(x,t) \in \Gamma \times (0,T), \\
[\beta(x)\nabla u \cdot n] &= g_2, &(x,t) \in \Gamma \times (0,T), \\
u(x,t) &= g(x,t), &(x,t) \in \partial\Omega \times (0,T), \\
u(x,0) &= u_0(x), &x \in \Omega
\end{align}

Unlike the elliptic case, the interface $\Gamma$ may be time-dependent, moving with a known motion. This can model scenarios like heat conduction in composite materials with dynamic interfaces.

\section{Local Randomized Neural Networks Methods}

\subsection{Randomized Neural Networks}
A fully connected neural network $\Psi: \mathbb{R}^{n_0} \rightarrow \mathbb{R}^{n_D}$ with depth $D$ is defined as:
\begin{align}
\Psi_0(x) &= x, \\
\Psi_l(x) &= \rho(W_l\Psi_{l-1} + b_l), \quad l = 1,\ldots,D-1, \\
\Psi &= \Psi_D = W_D\Psi_{D-1}
\end{align}

Here, $\rho$ is the activation function, and $W_l \in \mathbb{R}^{n_l \times n_{l-1}}$ and $b_l \in \mathbb{R}^{n_l}$ are the weights and biases in the $l$-th layer.

In a Randomized Neural Network (RNN), the weights and biases in the hidden layers are randomly generated and fixed, with only the weights of the output layer being adjustable. Any function $u_\rho \in \mathcal{N}_\rho^D$ (with one neuron in the output layer) can be expressed as a linear combination of basis functions:
\begin{equation}
u_\rho = \sum_{i=1}^{n_{D-1}} \alpha_i \psi_i
\end{equation}
where $\psi_i$ is the $i$-th component of the vector $\Psi_{D-1}$.

\subsection{LRNNs Method for Elliptic Interface Problem}

To solve the elliptic interface problem, the LRNNs method uses one RNN for each subdomain. For a domain split into two subdomains by interface $\Gamma$, two RNNs approximate the true solution in each subdomain:
\begin{equation}
u_\rho^1 = \sum_{k=1}^{m} \alpha_k^1 \psi_k^1, \quad u_\rho^2 = \sum_{k=1}^{m} \alpha_k^2 \psi_k^2
\end{equation}

Substituting these approximations into the elliptic interface problem and discretizing at randomly sampled collocation points ($N_1$ points in $\Omega$, $N_2$ points on $\Gamma$, and $N_3$ points on $\partial\Omega$), we obtain a linear system:
\begin{align}
AX &= F, \\
BX &= G_1, \\
CX &= G_2, \\
DX &= G_D
\end{align}

where $A, B, C, D$ are matrices of order $N_1 \times 2m$, $N_2 \times 2m$, $N_2 \times 2m$, and $N_3 \times 2m$, respectively, and $X = (\alpha_1^1, \ldots, \alpha_m^1, \alpha_1^2, \ldots, \alpha_m^2)^T$ is the unknown vector.

The derivatives are approximated using finite difference schemes. For example, the first and second-order partial derivatives with respect to $x$ are approximated as:
\begin{align}
\frac{\partial u_\rho(x,y)}{\partial x} &= \frac{u_\rho(x+h_x,y) - u_\rho(x-h_x,y)}{2h_x}, \\
\frac{\partial^2 u_\rho(x,y)}{\partial x^2} &= \frac{u_\rho(x+h_x,y) - 2u_\rho(x,y) + u_\rho(x-h_x,y)}{h_x^2}
\end{align}

The weights of the RNNs' output layers are obtained by solving a least-squares problem with this linear system.

To account for the importance of different conditions, weighting factors $\gamma_1$, $\gamma_2$, and $\gamma_3$ can be introduced:
\begin{align}
AX &= F, \\
\gamma_1 BX &= \gamma_1 G_1, \\
\gamma_2 CX &= \gamma_2 G_2, \\
\gamma_3 DX &= \gamma_3 G_D
\end{align}

These weights can be adjusted to emphasize the interface and boundary conditions, which are often crucial for accuracy.

\subsection{Mixed LRNNs Method}

By introducing another vector-valued function $p = \beta(x)\nabla u$, the elliptic interface problem can be rewritten in mixed form:
\begin{align}
-\nabla \cdot p &= f, &x \in \Omega, \\
p - \beta(x)\nabla u &= 0, &x \in \Omega, \\
[u] &= g_1, &x \in \Gamma, \\
[p \cdot n] &= g_2, &x \in \Gamma, \\
u &= g_D, &x \in \partial\Omega
\end{align}

Additional RNNs are needed to approximate $p$:
\begin{align}
u_\rho^1 &= \sum_{k=1}^{m} \alpha_k^1 \psi_k^1, \quad u_\rho^2 = \sum_{k=1}^{m} \alpha_k^2 \psi_k^2, \\
p_\rho^1 &= \left(\sum_{k=1}^{m} \tau_k^{1,1} \xi_k^{1,1}, \sum_{k=1}^{m} \tau_k^{1,2} \xi_k^{1,2}\right), \\
p_\rho^2 &= \left(\sum_{k=1}^{m} \tau_k^{2,1} \xi_k^{2,1}, \sum_{k=1}^{m} \tau_k^{2,2} \xi_k^{2,2}\right)
\end{align}

Substituting these into the mixed form and adding weights $\gamma_i$, we get a larger linear system with matrices of order $N_1 \times 6m$, $N_2 \times 6m$, and $N_3 \times 6m$.

\subsection{Space-time LRNNs Method for Parabolic Interface Problem}

For parabolic interface problems, the LRNNs method uses a space-time approach where time and space variables are treated equally. The interface becomes a surface in the space-time domain, dividing it into two subdomains. Two RNNs approximate the solution:
\begin{equation}
u_\rho^1 = \sum_{k=1}^{m} \alpha_k^1 \psi_k^1, \quad u_\rho^2 = \sum_{k=1}^{m} \alpha_k^2 \psi_k^2
\end{equation}

This leads to a linear system:
\begin{align}
AX &= F, \\
\gamma_1 BX &= \gamma_1 G_1, \\
\gamma_2 CX &= \gamma_2 G_2, \\
\gamma_3 DX &= \gamma_3 G_D, \\
\gamma_4 EX &= \gamma_4 U_0
\end{align}

where $A, B, C, D, E$ are matrices of order $N_1 \times 2m$, $N_2 \times 2m$, $N_2 \times 2m$, $\frac{4N_3}{5} \times 2m$, and $\frac{N_3}{5} \times 2m$, respectively.

This space-time approach avoids time-stepping iteration and accumulation errors, solving the entire problem in one least-squares solution.

\section{Numerical Examples and Results}

\subsection{Two-dimensional Flower Shape Interface Problem}
A domain $\Omega = (-1,1)^2$ with a flower-shaped interface:
\begin{equation}
\Gamma: (x-x_c)^2 + (y-y_c)^2 = r^2(\theta), \quad r(\theta) = 0.4 + 0.2\sin(20\theta), \quad \theta \in [0,2\pi)
\end{equation}

The exact solution is:
\begin{equation}
u(x) = 
\begin{cases}
\frac{1}{\beta_1}e^{xy}, & x \in \Omega_1, \\
\frac{1}{\beta_2}\sin(x)\sin(y), & x \in \Omega_2
\end{cases}
\end{equation}

Two RNNs with tanh activation functions approximated the solution. Results showed:
\begin{itemize}
    \item Error decreased with increasing neurons in hidden layers
    \item For $\beta_1 = 0.0001$ and $\beta_2 = 10000$, relative $L^2$ error reached $10^{-9}$
    \item Method performed well even with diffusion coefficients having large variations
    \item Mixed LRNNs method achieved even smaller errors when sufficient sampling points were used
\end{itemize}

\subsection{Three-dimensional Interface Problem}
A domain $\Omega = (-1,1)^3$ with a spherical interface:
\begin{equation}
\Gamma: x^2 + y^2 + z^2 = 0.75^2
\end{equation}

The exact solution:
\begin{equation}
u(x) = 
\begin{cases}
5e^{x^2+y^2+z^2} + 20, & x \in \Omega_1, \\
10(x+y+z), & x \in \Omega_2
\end{cases}
\end{equation}

Results showed:
\begin{itemize}
    \item For $\beta_1 = 1$ and $\beta_2 = 1000$, relative $L^2$ error reached $10^{-5}$
    \item Method was efficient, taking only about 6.8 seconds compared to 22.4 seconds for virtual element method
    \item Performance was robust across different diffusion coefficients
\end{itemize}

\subsection{Multiple Interfaces and High-dimensional Problems}
The method successfully handled:
\begin{itemize}
    \item Problems with three interfaces, using four RNNs with relative $L^2$ error reaching $10^{-9}$
    \item High-dimensional problems up to 20 dimensions with hyperplane interfaces
    \item Maintained accuracy even as dimension increased, with degrees of freedom not growing exponentially
\end{itemize}

\subsection{Time-dependent Interface Problems}
For parabolic problems with both fixed and moving interfaces, the space-time LRNNs approach showed:
\begin{itemize}
    \item High accuracy with no error accumulation over time
    \item Relative $L^2$ error reaching $10^{-7}$ for problems with fixed interfaces
    \item Effective handling of dynamic interfaces with evolving geometry
    \item Robustness to different diffusion coefficients
\end{itemize}

\section{Implications and Advantages}

\subsection{Computational Efficiency}
The LRNNs method offers significant efficiency advantages:
\begin{itemize}
    \item Avoids expensive optimization processes required by deep neural networks
    \item Solves linear least-squares problems instead of non-linear optimization
    \item Achieves high accuracy with fewer degrees of freedom
    \item Demonstrates faster computation times compared to traditional methods
\end{itemize}

\subsection{Handling Complex Geometries}
Being mesh-free, the method excels at problems with complex geometries:
\begin{itemize}
    \item Eliminates the need for complex interface meshing and fitting
    \item Handles flower-shaped, spherical, and multiple interfaces effectively
    \item Adapts easily to dynamic interfaces that change over time
    \item Maintains accuracy even with nearly touching or intersecting interfaces
\end{itemize}

\subsection{Solution Quality and Robustness}
The LRNNs approach demonstrates several advantages in solution quality:
\begin{itemize}
    \item Achieves spectral-like accuracy for smooth solutions
    \item Handles large variations in diffusion coefficients (ratios up to $10^{12}$)
    \item Space-time approach eliminates temporal error accumulation
    \item Performance improves predictably with increased sampling points and neurons
\end{itemize}

\subsection{Scalability to Higher Dimensions}
Unlike many traditional methods, LRNNs show good scalability:
\begin{itemize}
    \item Successfully solves problems up to 20 dimensions
    \item Degrees of freedom scale with solution complexity rather than exponentially with dimension
    \item Maintains reasonable accuracy even in high dimensions
    \item Sampling strategy adapts naturally to higher-dimensional problems
\end{itemize}

The LRNNs method represents a significant advancement in solving interface problems, combining the flexibility of neural networks with the efficiency of randomized algorithms, while avoiding the optimization challenges of traditional deep learning approaches.

\end{document} 